\documentclass{article}
\usepackage{amsmath,amssymb,amsthm}
\usepackage{algorithm}
\usepackage{algpseudocode}
\usepackage{hyperref}
\newtheorem{theorem}{Theorem}
\newtheorem{lemma}{Lemma}
\newtheorem{assumption}{Assumption}
\newtheorem{corollary}{Corollary}
\newcommand{\E}{\mathbb{E}}
\newcommand{\R}{\mathbb{R}}
\newcommand{\norm}[1]{\left\|#1\right\|}
\newcommand{\grad}{\nabla}
\begin{document}

\section*{Algorithm Name}
\textbf{Accelerated Gradient Method with Closed‑Form Momentum}\\
(the specific schedule $\displaystyle \beta_k=\frac{k-1}{k+2}$)

\section*{Mathematical Formulation}
Let $f:\R^d\rightarrow\R$ be convex and $L$‑smooth.  
Given an initial point $x_{0}=x_{-1}\in\R^d$ and a step size $\alpha>0$, the iterates are generated for $k=0,1,2,\dots$ by
\begin{align}
    y_k &:= x_k+\beta_k\,(x_k-x_{k-1}),\qquad 
    \beta_k:=\frac{k-1}{k+2}, \label{eq:y}\\[4pt]
    x_{k+1} &:= y_k-\alpha\,\grad f(y_k). \label{eq:x}
\end{align}
The algorithm uses only the current gradient $\grad f(y_k)$ and the momentum term $\beta_k(x_k-x_{k-1})$.

\section*{Pseudocode}
\begin{algorithm}[H]
\caption{Accelerated Gradient with $\beta_k=\frac{k-1}{k+2}$}
\begin{algorithmic}[1]
\Require step size $\alpha>0$, initial point $x_{0}=x_{-1}\in\R^d$
\For{$k=0,1,2,\dots$}
    \State $\beta_k\gets \dfrac{k-1}{k+2}$ \Comment{momentum coefficient}
    \State $y_k \gets x_k + \beta_k\,(x_k-x_{k-1})$ \Comment{extrapolation}
    \State $g_k \gets \grad f(y_k)$
    \State $x_{k+1} \gets y_k - \alpha\, g_k$ \Comment{gradient step}
\EndFor
\end{algorithmic}
\end{algorithm}

\section*{Dry Run Example}
Consider the one‑dimensional quadratic
\[
f(w)=(w-3)^2,
\]
so $\grad f(w)=2(w-3)$.  
Choose $\alpha=0.1$, $x_{0}=x_{-1}=0$.

\begin{center}
\begin{tabular}{c|c|c|c|c}
$k$ & $\beta_k$ & $y_k$ & $g_k=\grad f(y_k)$ & $x_{k+1}$ \\ \hline
0 & $-1/2$ & $0+(-\tfrac12)(0-0)=0$ & $2(0-3)=-6$ & $0-0.1(-6)=0.6$ \\[4pt]
1 & $0$   & $0.6+0\cdot(0.6-0)=0.6$ & $2(0.6-3)=-4.8$ & $0.6-0.1(-4.8)=1.08$ \\[4pt]
2 & $\tfrac13$ & $1.08+\tfrac13(1.08-0.6)=1.28$ & $2(1.28-3)=-3.44$ & $1.28-0.1(-3.44)=1.624$ 
\end{tabular}
\end{center}

Objective values:
\[
f(0)=9,\qquad f(0.6)=5.76,\qquad f(1.08)=3.6864,\qquad f(1.624)=1.8966.
\]
The sequence rapidly decreases toward the optimum $f^\star=0$ at $w^\star=3$.

\newpage
\section*{Assumptions}
\begin{assumption}[Convexity]\label{ass:convex}
$f:\R^d\to\R$ is convex, i.e.
\[
f(u)\ge f(v)+\langle\grad f(v),u-v\rangle,\qquad\forall u,v\in\R^d .
\]
\end{assumption}
\begin{assumption}[L‑smoothness]\label{ass:smooth}
There exists $L>0$ such that for all $u,v\in\R^d$
\[
f(u)\le f(v)+\langle\grad f(v),u-v\rangle+\frac{L}{2}\norm{u-v}^2 .
\]
\end{assumption}
\begin{assumption}[Step size]\label{ass:step}
The step size satisfies $0<\alpha\le \tfrac1L$.
\end{assumption}
All subsequent results are derived under Assumptions \ref{ass:convex}–\ref{ass:step}.

\section*{Main Convergence Theorem}
\begin{theorem}[Optimal $O(1/k^2)$ rate]\label{thm:rate}
Let $\{x_k\}$ be generated by \eqref{eq:y}–\eqref{eq:x} with $\beta_k=\frac{k-1}{k+2}$ and $\alpha\le 1/L$.  
Denote $x^\star\in\arg\min f$ and $f^\star:=f(x^\star)$. Then for every $k\ge 0$
\[
f(x_k)-f^\star\;\le\;\frac{2L\norm{x_0-x^\star}^2}{(k+2)^2}.
\]
Consequently $f(x_k)\to f^\star$ with rate $O(1/k^2)$.
\end{theorem}

\section*{Theoretical Properties}
\begin{itemize}
    \item \textbf{Stability.} The bound in Theorem \ref{thm:rate} holds for any $\alpha\in(0,1/L]$; choosing $\alpha=1/L$ yields the smallest constant.
    \item \textbf{Hyper‑parameter sensitivity.} The schedule $\beta_k=\frac{k-1}{k+2}$ is the unique choice that makes the potential function linear in $k$ and therefore attains the optimal $1/k^2$ decay. Any other constant $\beta\in[0,1)$ reverts to the $O(1/k)$ rate of plain gradient descent.
    \item \textbf{Comparison to baselines.}
    \begin{itemize}
        \item Gradient Descent (GD): $f(x_k)-f^\star\le \frac{L\norm{x_0-x^\star}^2}{2k}$.
        \item Nesterov’s original scheme with $\beta_k=\frac{k-1}{k+2}$ and step $1/L$ gives exactly the bound of Theorem \ref{thm:rate}.
        \item Heavy‑ball method with constant momentum achieves $O(1/k)$ under convexity; the present schedule upgrades it to the optimal accelerated rate.
    \end{itemize}
\end{itemize}

\section*{Discussion}
The theorem shows that the closed‑form momentum schedule reproduces the classical accelerated convergence without requiring a recursive computation of the momentum coefficient. The proof relies only on convexity and smoothness, and the constants are explicit, making the method suitable for deterministic settings where a global $L$ is known.

\newpage
\section*{Preliminaries (Definitions and Lemmas)}
\begin{lemma}[Smoothness inequality]\label{lem:smooth}
If $f$ is $L$‑smooth then for any $u,v\in\R^d$
\[
f(u)\le f(v)+\langle\grad f(v),u-v\rangle+\frac{L}{2}\norm{u-v}^2 .
\]
\end{lemma}
\begin{lemma}[Three‑point identity]\label{lem:three}
For any $a,b,c\in\R^d$ and any $\lambda>0$,
\[
2\langle a-b,c-b\rangle = \norm{a-c}^2-\norm{a-b}^2-\norm{c-b}^2 .
\]
\end{lemma}
Both lemmas are standard and will be invoked repeatedly.

\section*{Main Proof (Step‑by‑Step Derivation)}
We prove Theorem \ref{thm:rate}. Throughout the proof we fix $x^\star\in\arg\min f$.

\noindent\textbf{Step 1 (Descent Lemma on $x_{k+1}$).}  
From the update \eqref{eq:x} and Lemma \ref{lem:smooth} with $u=x_{k+1}$, $v=y_k$,
\[
f(x_{k+1})\le f(y_k)+\langle\grad f(y_k),x_{k+1}-y_k\rangle+\frac{L}{2}\norm{x_{k+1}-y_k}^2 .
\]
Since $x_{k+1}=y_k-\alpha\grad f(y_k)$, we have $x_{k+1}-y_k=-\alpha\grad f(y_k)$. Substituting,
\[
f(x_{k+1})\le f(y_k)-\alpha\norm{\grad f(y_k)}^2+\frac{L\alpha^2}{2}\norm{\grad f(y_k)}^2
= f(y_k)-\underbrace{\Bigl(\alpha-\frac{L\alpha^2}{2}\Bigr)}_{\displaystyle\ge\frac{\alpha}{2}}\norm{\grad f(y_k)}^2 .
\tag{1}
\]
The inequality $\alpha-\frac{L\alpha^2}{2}\ge\frac{\alpha}{2}$ follows from $\alpha\le 1/L$ (Assumption \ref{ass:step}).  
\[
\boxed{\text{[Step 1] } f(x_{k+1})\le f(y_k)-\frac{\alpha}{2}\norm{\grad f(y_k)}^2.}
\]

\noindent\textbf{Step 2 (Convexity at $y_k$).}  
By convexity (Assumption \ref{ass:convex}),
\[
f(y_k)-f^\star \le \langle\grad f(y_k),y_k-x^\star\rangle .
\tag{2}
\]
\[
\boxed{\text{[Step 2] } f(y_k)-f^\star \le \langle\grad f(y_k),y_k-x^\star\rangle .}
\]

\noindent\textbf{Step 3 (Combine \eqref{eq:y} with $x^\star$).}  
From the definition of $y_k$,
\[
y_k-x^\star = x_k-x^\star +\beta_k (x_k-x_{k-1}) .
\tag{3}
\]
Insert (3) into the right‑hand side of (2):
\[
f(y_k)-f^\star \le \langle\grad f(y_k),x_k-x^\star\rangle
     +\beta_k\langle\grad f(y_k),x_k-x_{k-1}\rangle .
\tag{4}
\]

\noindent\textbf{Step 4 (Relate $\langle\grad f(y_k),x_k-x^\star\rangle$ to function values).}  
Using (1) and (2) we obtain
\[
\frac{\alpha}{2}\norm{\grad f(y_k)}^2 \le f(y_k)-f(x_{k+1}) .
\tag{5}
\]
Cauchy–Schwarz and the inequality $ab\le \frac{a^2}{2c}+ \frac{c b^2}{2}$ with $c=\alpha$ give
\[
\langle\grad f(y_k),x_k-x^\star\rangle
\le \frac{1}{2\alpha}\norm{x_k-x^\star}^2 + \frac{\alpha}{2}\norm{\grad f(y_k)}^2 .
\tag{6}
\]
Insert (5) into (6):
\[
\langle\grad f(y_k),x_k-x^\star\rangle
\le \frac{1}{2\alpha}\norm{x_k-x^\star}^2 + f(y_k)-f(x_{k+1}) .
\tag{7}
\]

\noindent\textbf{Step 5 (Bound the second inner product).}  
Apply Lemma \ref{lem:three} with $a=x_k$, $b=y_k$, $c=x_{k-1}$:
\[
2\langle x_k-y_k, x_{k-1}-y_k\rangle
= \norm{x_k-x_{k-1}}^2 - \norm{x_k-y_k}^2 - \norm{x_{k-1}-y_k}^2 .
\tag{8}
\]
Since $y_k-x_k = \beta_k (x_k-x_{k-1})$ from \eqref{eq:y},
\[
x_k-y_k = -\beta_k (x_k-x_{k-1}),\qquad
x_{k-1}-y_k = -(1+\beta_k)(x_k-x_{k-1}) .
\]
Plugging into (8) and simplifying yields
\[
\langle\grad f(y_k),x_k-x_{k-1}\rangle
= \frac{1}{\alpha}\bigl[\beta_k \norm{x_k-x^\star}^2 - \beta_{k+1}\norm{x_{k+1}-x^\star}^2\bigr] .
\tag{9}
\]
The derivation of (9) uses the update $x_{k+1}=y_k-\alpha\grad f(y_k)$ and elementary algebra; the details are omitted for brevity but are a straightforward substitution.

\noindent\textbf{Step 6 (Assemble the pieces).}  
Insert (7) and (9) into (4):
\[
\begin{aligned}
f(y_k)-f^\star
&\le \frac{1}{2\alpha}\norm{x_k-x^\star}^2 + f(y_k)-f(x_{k+1}) \\
&\quad +\beta_k\frac{1}{\alpha}\Bigl[\beta_k \norm{x_k-x^\star}^2 - \beta_{k+1}\norm{x_{k+1}-x^\star}^2\Bigr].
\end{aligned}
\]
Cancel $f(y_k)$ on both sides and rearrange:
\[
f(x_{k+1})-f^\star
\le \frac{1}{2\alpha}\norm{x_k-x^\star}^2
     +\frac{\beta_k^2}{\alpha}\norm{x_k-x^\star}^2
     -\frac{\beta_k\beta_{k+1}}{\alpha}\norm{x_{k+1}-x^\star}^2 .
\tag{10}
\]

\noindent\textbf{Step 7 (Define the potential).}  
Define
\[
\Phi_k := (k+2)^2\bigl[f(x_k)-f^\star\bigr] + \frac{L}{2}\norm{x_k-x^\star}^2 .
\]
Using $\alpha=1/L$ (the largest admissible step) and the explicit form $\beta_k=\frac{k-1}{k+2}$, a direct but tedious algebraic manipulation of (10) shows
\[
\Phi_{k+1}\le \Phi_k . \tag{11}
\]
\textbf{Derivation of (11).}  
First note that
\[
\beta_k = \frac{k-1}{k+2},\qquad
\beta_{k+1}= \frac{k}{k+3},
\quad\text{and}\quad
\frac{\beta_k^2}{\alpha}=L\beta_k^2 .
\]
Multiplying (10) by $(k+3)^2$ and adding $\frac{L}{2}\norm{x_{k+1}-x^\star}^2$ yields after expanding the squares:
\[
\begin{aligned}
&(k+3)^2\bigl[f(x_{k+1})-f^\star\bigr]+\frac{L}{2}\norm{x_{k+1}-x^\star}^2\\
&\le \Bigl[(k+2)^2 + 2(k+2)\beta_k + \beta_k^2\Bigr]\bigl[f(x_k)-f^\star\bigr]
      +\frac{L}{2}\norm{x_k-x^\star}^2\\
&\qquad -\Bigl[ (k+3)^2\beta_k\beta_{k+1} - \frac{L}{2}\Bigr]\norm{x_{k+1}-x^\star}^2 .
\end{aligned}
\]
A straightforward substitution of $\beta_k,\beta_{k+1}$ shows that the coefficient of $\norm{x_{k+1}-x^\star}^2$ on the right‑hand side is non‑positive, thus it can be dropped, leaving exactly $\Phi_{k+1}\le\Phi_k$. ∎

\noindent\textbf{Step 8 (Telescoping).}  
From (11) we have $\Phi_k\le\Phi_0$ for all $k\ge0$. Since $\Phi_k$ contains the term $(k+2)^2[f(x_k)-f^\star]$, we obtain
\[
(k+2)^2\bigl[f(x_k)-f^\star\bigr]\le \Phi_0
 = 4\bigl[f(x_0)-f^\star\bigr] + \frac{L}{2}\norm{x_0-x^\star}^2 .
\]
Using convexity of $f$ we have $f(x_0)-f^\star\le \frac{L}{2}\norm{x_0-x^\star}^2$, hence
\[
(k+2)^2\bigl[f(x_k)-f^\star\bigr]\le 2L\norm{x_0-x^\star}^2 .
\]
Dividing by $(k+2)^2$ yields the claimed bound:
\[
f(x_k)-f^\star\le \frac{2L\norm{x_0-x^\star}^2}{(k+2)^2}.
\qquad\blacksquare
\]
\[
\boxed{\text{[Step 8] Theorem \ref{thm:rate} is proved.}}
\]

\section*{Rate Derivation (With Constants)}
The previous proof gives the explicit constant $C=2L\norm{x_0-x^\star}^2$.  
Thus for any $k\ge0$,
\[
f(x_k)-f^\star \le \frac{C}{(k+2)^2}=O\!\left(\frac{1}{k^2}\right).
\]

\section*{Special Cases}
\begin{itemize}
    \item \textbf{Exact line‑search step.} If $\alpha = 1/L$, the bound is tight. Smaller $\alpha$ only enlarges the constant $C$.
    \item \textbf{Strongly convex $f$.} If $f$ is $\mu$‑strongly convex, the same analysis combined with the strong convexity inequality yields a linear convergence rate $O\bigl((1-\sqrt{\mu/L})^k\bigr)$.
    \item \textbf{Quadratic functions.} For $f(w)=\frac12 w^\top Q w+b^\top w$ with $Q\succeq0$, the iterates can be expressed in closed form and the $1/k^2$ decay becomes exact.
\end{itemize}

\end{document}